% Section 1: introduction (short-paper outline, RESEARCH_RECORD.md). Statements about the literature mean
% "not found in the texts read"; gate S1 was closed with a stated risk (two sources unread), see Sec. 7.

\section{Introduction}\label{sec:intro}

A posteriori error estimation plays an important role in assessing the reliability of numerical approximations to contact and obstacle problems.
In unilateral contact mechanics, an approximate displacement must satisfy non-penetration constraints, while the associated contact forces must
remain nonnegative and obey complementarity conditions. These inequalities distinguish contact problems from unconstrained elliptic equations and
introduce additional difficulties in error estimation. Established error certificates address these difficulties by combining an equilibrium
residual with a complementarity contribution involving a feasible displacement approximation and a nonnegative multiplier. Such estimates have
been developed in the functional framework for obstacle problems \cite{repin2000,repin2007}, through finite element error estimators
\cite{veeser2001,braess2005,krause2015}, and for discrete variational inequalities and reduced-order contact models
\cite{zhang2016cr,zhang2016thesis}. Since the multiplier used in the certificate need not coincide with that of the numerical approximation, it
can be selected to improve the resulting bound. One established strategy is to minimise a convex quadratic surrogate of the certificate, leading
to a linear complementarity problem \cite{repin2007,zhang2016thesis}. We refer to this strategy as S1.

The reliability and efficiency of a posteriori error estimators for obstacle and Signorini problems have been investigated extensively. Existing
finite element analyses establish efficiency estimates involving primal and multiplier errors and identify the influence of contact regions on
estimator behaviour \cite{veeser2001,braess2005,krause2015}. Reduced-basis methods for variational inequalities have also revealed a possible
square-root difference between error bounds \cite{haasdonk2012}. A duality approach with strictly feasible approximations reports bounds that numerically follow the convergence rate of the reduced-basis error in two model problems, without deriving an efficiency bound \cite{zhang2014arxiv}. More generally, primal--dual gap formulations and error identities provide rigorous
relationships among equilibrium residuals, approximation errors and contact-complementarity terms \cite{bartels2025,repin2018}. These results
establish the mathematical foundations of contact-aware error estimation. They also show that the interaction between displacement accuracy and
contact consistency is an important consideration when interpreting an error bound.

Despite these developments, the efficiency of the particular feasible-pair certificate obtained through convex-surrogate multiplier optimisation
remains less explicitly characterised. To the best of our knowledge, the literature does not provide an explicit characterisation of when the
optimised certificate coincides with the true displacement error. Furthermore, it neither demonstrates that convergence-rate degradation may
persist under unrestricted multiplier optimisation nor analyses the asymptotic certificate floor caused by omitted strictly active contact-force
components. While established primal--dual energy identities provide closely related estimates, they leave open key questions regarding the
surrogate-minimising multiplier.

To address this gap, the present study investigates the relationship between the optimised certificate, the true displacement error and the
complementarity between exact contact forces and approximate contact gaps. We first establish an efficiency estimate that quantifies this relationship
for a general symmetric positive-definite inner product. For the stiffness-based inner product, we interpret the surrogate-minimising multiplier as
the contact-force solution of a problem with a shifted obstacle. This interpretation enables a precise characterisation of the conditions under
which the optimised multiplier recovers the exact multiplier and the resulting certificate equals the true energy error. We then examine the
limitations of multiplier optimisation through an explicit counterexample demonstrating unavoidable worst-case square-root degradation within the
considered certificate family. A separate analysis shows that restricting the multiplier support so that a strictly active contact constraint is
omitted can produce a positive asymptotic certificate floor. These findings are complemented by bounds that quantify the contribution of missed
contact constraints to certificate efficiency.

The theoretical results are examined numerically using two finite element contact benchmarks, controlled approximation sequences and randomly
generated constrained quadratic problems. We also consider elementary feasibility and snap-to-contact corrections of reduced-order predictions to
illustrate the relationship between active-set consistency, certificate efficiency and displacement accuracy. The investigation is analytical in
scope: it does not propose a new error certificate, multiplier-selection algorithm or reduced-order method. Moreover, the conditions for exact
certification involve quantities from the unknown full-order solution and are therefore intended to explain certificate behaviour rather than
provide a directly applicable online exactness test. The numerical experiments are designed to verify the derived properties rather than demonstrate
computational acceleration.

The remainder of the paper is organised as follows. Section~\ref{sec:setting} presents the mathematical formulation and feasible-pair error
certificate. Section~\ref{sec:analysis} investigates certificate efficiency and conditions for exact certification. Section~\ref{sec:failure}
examines square-root convergence degradation, missed contact constraints and the effects of restricted multiplier support.
Section~\ref{sec:results} provides numerical verification of the theoretical results, while Section~\ref{sec:application} illustrates their
application to corrections of reduced-order contact predictions. Section~\ref{sec:conclusions} summarises the principal findings, limitations
and directions for future research. Appendix~\ref{app:e3} presents the random-instance verification results and supplementary numerical diagnostics.
